% ============================================================
% Page 6 — Alban aux « siècles bénédictins »
% Transcription réalisée à partir du scan fourni.
% Texte conservé dans sa forme originale ; aucune correction
% grammaticale ou reformulation n'a été effectuée.
% ============================================================

\thispagestyle{plain}

\vspace*{12mm}

\begin{center}
    {\Large\bfseries Alban aux « siècles bénédictins »}
\end{center}

\vspace{18mm}

\noindent
Un bout de parchemin trouvé dans un coin d'un immeuble, aujourd'hui
propriété de M. Délécouls, démontre que, dans la première moitié du
XIII\textsuperscript{ème} siècle, on vénérait ici une relique de
St Vincent, patron des vignerons. Le corps de ce martyr se trouvais
à l'Abbaye de Castres. Selon toute vraisemblance, ce sont les moines
venus de ce monastère qui en apportèrent une relique constituée par
un lambeau de tissu prélevé sur le vêtement de St Vincent. C'est ce
dont fait foi le bout de parchemin en question qu'on peut traduire
« de la tunique du bienheureux Vincent et du tombeau du Seigneur et
du tombeau de N.D. ».

\par\vspace{12mm}

\noindent
Est-ce à dire que ce sont ces bénédictins qui implantèrent ici la
vigne -- comme ce fut le cas pour le Gaillac -- le Champagne ?
Ce n'est pas impossible. Il est vrai que les vignobles sont rares
aux abords d'Alban. Le 12 octobre 1537, noble François Caudières
d'Alban vendait à Bernard Olim de la Mouline à fer une vigne à
Malagousse (parages de l'actuel « Parrot » (archives notariales
d'Albi).

\vfill
